% =============================================================================
% threat_model.tex -- Threat model.
%
% REVISION CONTRACT (advisor feedback, 2026-07):
% * Every term used here is defined in Section 2.4 (sec:concepts) and referenced, not
%   re-introduced. No new technical vocabulary is created in this section.
% * The two parties (auditor and adversary) are separated explicitly, because they have
%   different goals and different access, and the paper's instrument-validity question
%   is really about the auditor.
% * Success criteria are tied to the assumption chain of Section 2.5.
% =============================================================================

\section{Threat Model}
\label{sec:threat}

This section states who is acting, what they know, what they are trying to achieve, and
what counts as success. All technical terms used here are defined in
Section~\ref{sec:concepts}.

\subsection{Two parties with different problems}
Discussions of contamination often blur two roles that need separating, because they have
different access and different objectives.

The \textbf{auditor} is a defender. This is a model developer, an evaluation
organization, or an external reviewer who wants to know whether a benchmark leaked into
training and, more usefully, \emph{which} specific items a model has retained. The
auditor is willing to run cheap tests over many items and cannot afford to train
reference models. The auditor is the party this paper is really about, because the
question of whether a detector is a valid instrument is the auditor's question.

The \textbf{adversary} is an attacker who wants to recover content from the training
data. The adversary needs only the deployed model and does not need the training corpus.

A key asymmetry is that the auditor is trying to \emph{predict} what the adversary could
extract, without performing the extraction on every item. That is precisely why detector
validity matters, and it is the practical stake behind assumption A1 in
Section~\ref{sec:chain}.

\subsection{Goals, in increasing order of harm}
\begin{itemize}
 \item \textbf{G1: membership inference.} Decide whether one specific sequence, such as a
 benchmark item, document, or record, was in the training corpus.
 \item \textbf{G2: benchmark-level contamination confirmation.} Decide, with a controlled
 false-positive rate, whether an entire benchmark was trained on.
 \item \textbf{G3: extraction.} Cause the model to reproduce content that was in the
 training data. On a corpus with known sensitive content this is where privacy harm would
 occur.
\end{itemize}

G3 is the outcome that matters. G1 and G2 are of interest mainly insofar as they predict
G3, and whether they do is exactly what this paper measures. Under the terminology of
Section~\ref{sec:chain}, a detector that achieves G1 but fails to predict G3 is a
detector for which A1 does not hold in the sense the privacy framing requires.

\subsection{Knowledge and access}
We grade the detectors by the least access each one needs, using the definitions from
Section~\ref{sec:concepts}:

\begin{itemize}
 \item \textbf{Black-box.} Text in, text out, with no probabilities. Guided prompting
 falls here. We do not evaluate black-box detectors.
 \item \textbf{Gray-box.} The model's per-token log-probabilities for a supplied text.
 LOSS, Min-K\%, and the zlib ratio need only this.
 \item \textbf{White-box.} The full next-token probability distribution at each position.
 Min-K\%++ requires this, because it rescales each token against the alternatives the
 model considered.
\end{itemize}

The $n$-gram overlap test is different in kind, because it inspects the training corpus
rather than the model. We can run it only because The Pile is public, and we use it to
construct contamination labels rather than as an attack an adversary could mount.
Extraction (G3) requires only the ability to prompt the model and read its output.

For our experiments the auditor additionally knows the training corpus, which is what
makes membership ground truth rather than an inference. We treat this as a deliberate
methodological advantage rather than a realistic assumption, and Section~\ref{sec:eval}
explains why studying a model whose corpus is public is the only way to test A1 without
confounding.

\subsection{Success criteria}
\begin{itemize}
 \item \textbf{G1.} True-positive rate at $0.1\%$ and at $1\%$ false-positive rate, read
 from a logarithmic-axis ROC curve, with AUC reported as secondary and with bootstrap
 confidence intervals. Section~\ref{sec:concepts} explains why the low-false-positive
 operating point is the meaningful one for a privacy attack.
 \item \textbf{G2.} A permutation-test $p$-value below a pre-specified threshold, with a
 controlled false-positive rate~\citep{oren2024proving}.
 \item \textbf{G3.} A non-zero extraction rate, and then the paper's central criterion: a
 \emph{positive} association between a per-item detector score and the per-item extraction
 outcome that \emph{survives holding the model's loss fixed}, in the partial-correlation
 sense of Section~\ref{sec:concepts}.
\end{itemize}

The last criterion is the one that distinguishes a detector that genuinely identifies
which items are at risk from one that merely restates the model's loss in different
units. We fixed this criterion in advance of running the analysis, for the reason given
under pre-registration in Section~\ref{sec:concepts}.

\subsection{Out of scope}
We do not attack deployed production models for real third-party personal data. We do not
train or fine-tune models, and we propose no new detector. We do not claim to demonstrate
leakage of sensitive information, since assumptions A3 and A4 of
Section~\ref{sec:chain} are untested here. Differential privacy is discussed as the
producer-side mitigation this threat model motivates (Section~\ref{sec:dp}) and is not
implemented or evaluated.
